\chapter*{Список использованных сокращений}
\addcontentsline{toc}{chapter}{Список использованных сокращений}
\begin{description}
\item[BANC] Brain and Nerve Cord; реконструкция мозга и вентральной нервной цепочки взрослой самки Drosophila.
\item[ECAP] Вызванный составной потенциал действия (evoked compound action potential).
\item[GABA] $\gamma$-Аминомасляная кислота (gamma-aminobutyric acid).
\item[MANC] Male Adult Nerve Cord; реконструкция вентральной нервной цепочки взрослого самца Drosophila.
\item[RMSE] Среднеквадратическая ошибка (root mean square error).
\item[SCS] Стимуляция спинного мозга (spinal cord stimulation).
\end{description}

\chapter*{Список терминов}
\addcontentsline{toc}{chapter}{Список терминов}
\noindent динамическое представление: Формализованное описание выбранных временных характеристик нейронной активности.\par
\noindent коннектом: Совокупность сведений о нейронах и связях между ними.\par
\noindent кривая роста ECAP: Зависимость амплитуды вызванного составного потенциала действия от амплитуды стимуляции при фиксированных остальных параметрах.\par
\noindent оператор наблюдения: Преобразование состояния модели в величину, сопоставимую с измерением.\par
\noindent рекрутирование волокон: Вовлечение нервных волокон в вызванный стимуляцией ответ.\par
